\chapter{Acute Respiratory Distress Syndrome}
\index{Acute Respiratory Distress Syndrome}
\label{sec:Acute Respiratory Distress Syndrome}

\section{Overview}
Acute respiratory distress syndrome (ARDS) is a life-threatening form of acute lung injury characterized by diffuse alveolar damage, increased pulmonary vascular permeability, and severe hypoxemia. It typically develops within 6--72 hours of a direct or indirect lung insult and requires mechanical ventilation with positive end-expiratory pressure. ARDS carries a high mortality rate of 30--40\% and is a major cause of respiratory failure in critically ill patients worldwide.
\section{Classification}

\begin{description}[font=\normalfont\bfseries,leftmargin=3em,labelwidth=2.5em]
  \item[Cause] Direct or indirect lung injury triggering systemic inflammation
  \item[Organ System] Lungs/Respiratory
  \item[Pathophysiology] Diffuse alveolar damage from inflammatory cascade leading to increased capillary permeability, surfactant dysfunction, alveolar flooding, and hyaline membrane formation
  \item[Duration] Acute (days to weeks)
  \item[Onset] Sudden (within 6--72 hours of inciting event)
  \item[Mode of Transmission] Non-communicable
  \item[Clinical Course] Rapid onset of respiratory failure requiring intensive care
  \item[Severity] Severe (mild/moderate/severe based on PaO2/FiO2 ratio)
  \item[Extent of Disease] Diffuse (bilateral lung involvement)
  \item[Age Group] All ages
  \item[Epidemiology] Incidence approximately 10\% of ICU admissions; mortality 30--40\%
  \item[ICD-11 Code] CB00.0
\end{description}

\section{Causes}
Direct lung injuries causing ARDS include pneumonia (bacterial, viral, or fungal), aspiration of gastric contents, pulmonary contusion, near-drowning, and inhalation of toxic gases or smoke. Indirect injuries include sepsis (most common non-pulmonary cause), severe trauma with multiple transfusions, acute pancreatitis, burns, and drug overdoses. The common pathway involves activation of the innate immune system with neutrophil infiltration, cytokine release, and disruption of the alveolar-capillary barrier.

\section{Risk Factors}

\begin{itemize}[noitemsep]
  \item Sepsis (especially gram-negative bacteremia)
  \item Severe pneumonia (Streptococcus pneumoniae, influenza, SARS-CoV-2)
  \item Aspiration of gastric contents
  \item Major trauma or burns
  \item Acute pancreatitis
  \item Massive blood transfusion (transfusion-related acute lung injury, TRALI)
  \item Drug overdose (opioids, salicylates, tricyclic antidepressants)
  \item Near-drowning or smoke inhalation
  \item Severe acute pancreatitis
\end{itemize}

\section{Signs and Symptoms}

\subsection{Early symptoms}
\begin{itemize}[noitemsep]
  \item Early manifestations of acute respiratory distress syndrome may be subtle and nonspecific
\end{itemize}

\subsection{Common symptoms}
\begin{itemize}[noitemsep]
  \item \subsection{Common symptoms}
  \item \textbf{Common symptoms:}
  \item Severe shortness of breath (dyspnea) with rapid onset
  \item Tachypnea and labored breathing
  \item Hypoxemia refractory to supplemental oxygen
  \item Rapid, shallow respirations
  \item Confusion or altered mental status from cerebral hypoxia
  \item Cyanosis (bluish discoloration of skin and mucous membranes)
  \item Crackles (rales) on lung auscultation
  \item Profound fatigue and respiratory muscle weakness
  \item Pain or discomfort in the affected area
  \item Difficulty performing daily activities
\end{itemize}

\subsection{Less common symptoms}
\begin{itemize}[noitemsep]
  \item Atypical or less common presentations of acute respiratory distress syndrome may occur
\end{itemize}

\subsection{Emergency warning signs}
\begin{itemize}[noitemsep]
  \item Severe or worsening symptoms of acute respiratory distress syndrome require immediate medical evaluation
\end{itemize}
\section{Progression and Stages}
ARDS progresses through three overlapping phases: the exudative phase (days 1--7) characterized by alveolar flooding with protein-rich fluid, hyaline membrane formation, and neutrophilic inflammation; the proliferative phase (days 7--21) with fibroblast proliferation and organization of exudates; and the fibrotic phase (after day 21) with collagen deposition and potential for irreversible pulmonary fibrosis. Severity is graded by the Berlin definition based on the PaO2/FiO2 ratio with PEEP $\ge$5 cmH2O: mild (200--300 mmHg), moderate (100--200 mmHg), and severe ($<$100 mmHg).

\section{Complications}

\begin{itemize}[noitemsep]
  \item Ventilator-associated pneumonia
  \item Barotrauma and volutrauma (pneumothorax, pneumomediastinum)
  \item Multi-organ failure (renal, hepatic, cardiovascular)
  \item Pulmonary fibrosis in prolonged or severe cases
  \item Critical illness myopathy and polyneuropathy
  \item Post-intensive care syndrome (cognitive, psychiatric, and physical impairments)
  \item Reduced quality of life
  \item Emotional and psychological effects
\end{itemize}

\section{Diagnosis}

\begin{itemize}[noitemsep]
  \item Clinical diagnosis by Berlin criteria: acute onset ($<$1 week), bilateral opacities on chest imaging, respiratory failure not fully explained by cardiac failure or fluid overload
  \item Chest radiograph or CT scan showing bilateral diffuse infiltrates
  \item Arterial blood gas analysis measuring PaO2/FiO2 ratio for severity classification
  \item Echocardiogram to exclude cardiogenic pulmonary edema
  \item Bronchoalveolar lavage when infection or hemorrhage is suspected
  \item Lab tests including complete blood count, inflammatory markers, and cultures
  \item Imaging tests to monitor progression and complications
\end{itemize}

\section{Prevention}

\begin{itemize}[noitemsep]
  \item Early recognition and treatment of sepsis with appropriate antibiotics
  \item Lung-protective ventilation strategies in intubated patients
  \item Judicious fluid management and avoidance of fluid overload
  \item Early mobilization and oral care to reduce ventilator-associated complications
  \item Vaccination against respiratory pathogens (influenza, pneumococcus, SARS-CoV-2)
  \item Go for regular check-ups
  \item Follow your doctor's screening recommendations
\end{itemize}

\section{Treatment Options}

\begin{itemize}[noitemsep]
  \item Lung-protective mechanical ventilation with low tidal volumes (6 mL/kg ideal body weight)
  \item Positive end-expiratory pressure (PEEP) to maintain alveolar recruitment
  \item Prone positioning for severe ARDS to improve ventilation-perfusion matching
  \item Neuromuscular blockade (cisatracurium) in early severe ARDS
  \item Fluid-conservative strategy after initial resuscitation
  \item Corticosteroids (dexamethasone or methylprednisolone) in moderate-to-severe ARDS
  \item Extracorporeal membrane oxygenation (ECMO) for refractory hypoxemia
  \item Inhaled pulmonary vasodilators (nitric oxide, prostacyclin) as rescue therapy
  \item Treat underlying cause (antibiotics for infection, source control for sepsis)
  \item Lifestyle changes and self-care
  \item Surgery when needed
\end{itemize}

\section{Living with It}

\begin{itemize}[noitemsep]
  \item Follow your treatment plan as prescribed
  \item Keep up with regular appointments
  \item Adjust your daily habits as needed
  \item Reach out to family, friends, or support groups
  \item Keep learning about your condition
\end{itemize}

\section{Outlook}
Despite advances in critical care, ARDS carries a mortality of 30--40\% and survivors often face long-term physical, cognitive, and psychiatric sequelae. Among those who survive, pulmonary function typically improves over 6--12 months, though residual restrictive or diffusion defects and impaired quality of life may persist.
